\section{Introduction}
\label{sec:intro}

When asked to sort a list of countries alphabetically, a large language model succeeds effortlessly.
When asked to sort cities by elevation, the same model falters.
This asymmetry is striking: both tasks require the model to access knowledge and produce an ordered list, yet one is trivially easy and the other surprisingly hard.
What determines which properties an LLM can reliably order by?

This question has practical consequences.
LLMs are increasingly embedded in data-processing pipelines that involve sorting and ranking~\cite{patel2025semantic,zhao2025access}.
Knowing which orderings to trust---and which require external tools---is essential for building reliable systems.
The question also has scientific value: ordering accuracy provides a direct window into what factual knowledge LLMs have internalized and how precisely they can access it.

Prior work has examined LLM sorting ability, but only along narrow axes.
\citet{herbold2025sortbench} benchmark syntactic sorting of numbers and strings, finding that frontier models handle it well but struggle with adversarial cases like lexicographic sorting of number words.
\citet{zhao2025access} compare sorting algorithms (pointwise, pairwise, listwise) on a handful of factual properties such as NBA player heights and country populations, showing that pointwise methods work well when facts reside in parametric memory.
However, no existing work systematically tests ordering accuracy across a \emph{taxonomy} of property types to identify which features LLMs have easiest access to.

We address this gap with \benchname, a benchmark of 22 ordering tasks organized into five categories (syntactic, well-known factual, general knowledge, specific factual, and temporal) and three difficulty levels.
Each task presents a shuffled list of 6--10 items and asks the model to sort them by a specified property in a single pass.
We evaluate \gptlarge and \gptsmall on all tasks using three random shuffles per task and measure ordering quality with \kendalltau, \spearmanrho, and exact-match rate.

Our results reveal a clear \emph{property-accessibility hierarchy}.
\gptlarge achieves $\tau = 1.0$ on algorithmic operations (alphabetical sorting) and well-known sequences (planet order, Mohs hardness, historical chronologies), but drops to $\tau \approx 0.81$--$0.88$ on tasks requiring precise numeric recall (city elevations, river lengths, building heights).
\gptlarge significantly outperforms \gptsmall (mean $\tau = 0.948$ vs.\ $0.867$, $p = 0.008$), yet both models agree on which properties are hardest ($\rho = 0.61$, $p = 0.003$), indicating that the hierarchy reflects genuine differences in knowledge encoding rather than model-specific artifacts.

Our contributions are:
\begin{itemize}[leftmargin=*,itemsep=0pt,topsep=0pt]
    \item We introduce \benchname, a benchmark of 22 diverse ordering tasks spanning five property categories, designed to probe which features LLMs can most reliably access from parametric memory.
    \item We identify a property-accessibility hierarchy---from algorithmic operations ($\tau = 1.0$) through memorized sequences ($\tau = 1.0$) and approximate magnitudes ($\tau \approx 0.93$) down to precise numeric facts ($\tau \approx 0.83$)---that is consistent across models.
    \item We show that the key distinction is not whether knowledge is ``well-known'' versus ``obscure,'' but whether it is stored as a \emph{sequence} versus \emph{independent numeric values}, reframing how we think about factual accessibility in LLMs.
\end{itemize}
